\chapter{Proposta}
\label{sec:proposta}

% \section{Comparação com o Trabalho Proposto}

% Os estudos selecionados apresentam contribuições relevantes para a compreensão das dificuldades e necessidades de diferentes perfis de usuários durante a interação com sistemas digitais. Entretanto, as pesquisas analisadas estão inseridas em diferentes contextos de aplicação e abordam aspectos específicos do problema.

% Por exemplo, \cite{weighted_heuristic} concentra-se na avaliação quantitativa da usabilidade de websites para usuários idosos, relacionando heurísticas a medidas de desempenho. Já \cite{health_problem_solving} investiga a realização de tarefas complexas em um website governamental, enquanto \cite{social_media_elderly} aborda requisitos de interação em plataformas de mídia social.

% Outros estudos concentram-se em estratégias específicas de apoio ao usuário. \cite{kabassi2004personalised} apresentam um ambiente de treinamento adaptativo para adultos novatos, enquanto \cite{cowan2011exploring} investigam a influência do treinamento sobre a usabilidade percebida e a ansiedade de usuários novatos.

% Também são encontrados trabalhos que aplicam processos de desenvolvimento centrados no usuário e avaliações com participantes do público-alvo para construir soluções acessíveis e utilizáveis, como aquele feito por \cite{breaking_barriers}.

% Dessa forma, os estudos selecionados apresentam diferentes métricas, diretrizes, métodos de avaliação e estratégias de desenvolvimento, mas essas contribuições estão distribuídas entre diferentes contextos e objetivos de pesquisa.

% Nesse contexto, o presente trabalho propõe organizar essas contribuições em um guia de desenvolvimento direcionado a softwares e sites destinados a pessoas com baixa alfabetização informática. A proposta busca reunir métricas, recomendações de projeto e componentes de interface identificados na literatura, de maneira a transformar os conhecimentos encontrados nos estudos selecionados em uma referência de apoio ao desenvolvimento.